\textbf{Article information}

\textbf{Article title}

Materials Science Optimization Benchmark Dataset for High-dimensional,
Multi-objective, Multi-fidelity Optimization of CrabNet Hyperparameters

\textbf{Authors}

Sterling G. Baird\textsuperscript{1}*, Jeet N.
Parikh\textsuperscript{2}, Taylor D. Sparks\textsuperscript{1,3}

\textbf{Affiliations}

1. Materials Science \& Engineering, University of Utah, 122 S. Central
Campus Drive, \#304 Salt Lake City, Utah 84112

2. Northwood High School, 4515 Portola Pkwy, Irvine, CA 92620

3. Chemistry Department, University of Liverpool, Liverpool, L7 3NY,
United Kingdom

\textbf{Corresponding author's email address and Twitter handle}

\href{mailto:sterling.baird@utah.edu}{\nolinkurl{sterling.baird@utah.edu}}

@SterlingBaird1

\textbf{Keywords}

adaptive design, Bayesian optimization, formulation optimization,
PseudoCrab

\textbf{Abstract}

Benchmarks are crucial for driving progress in scientific disciplines.
To be effective, benchmarks should closely mimic real-world tasks while
being computationally efficient, allowing for accessibility and
repeatability. Developing surrogate models that can be indistinguishable
from the ground truth observation within the explored dataset bounds
dramatically reduces the computational burden of running benchmarks
without sacrificing quality, but this requires a large amount of initial
data. In the fields of materials science and chemistry, relevant
optimization tasks can be challenging due to their complexity, which
includes hierarchical, noisy, multi-fidelity, multi-objective,
high-dimensional, and non-linearly correlated variables. Additionally,
they may include mixed numerical and categorical variables that are
subject to linear and non-linear constraints. Simulating or
experimentally verifying such tasks can be difficult, which is why
benchmarks are essential. This study aimed to overcome these challenges
by generating 173219 quasi-random hyperparameter combinations across 23
hyperparameters and using them to train CrabNet on the Matbench
experimental band gap dataset (Computational runtime: 387 RTX-2080-Ti
GPU days). The results were stored in a free-tier shared MongoDB Atlas
dataset, creating a regression dataset that maps hyperparameter
combinations to metrics such as MAE, RMSE, computational runtime, and
model size for the CrabNet model trained on the Matbench experimental
band gap benchmark task. To simulate the actual simulations,
heteroskedastic noise was incorporated into the regression dataset, and
bad hyperparameter combinations were excluded. Percentile ranks were
computed within each group of identical parameter sets to capture
heteroskedastic noise, rather than assuming Gaussian noise as is done in
traditional approaches. This approach can be applied to other benchmark
datasets, bridging the gap between optimization benchmarks with low
computational overhead and realistically complex, real-world
optimization scenarios.

\textbf{Specifications table}

\begin{longtable}[]{@{}
  >{\raggedright\arraybackslash}p{(\columnwidth - 2\tabcolsep) * \real{0.2690}}
  >{\raggedright\arraybackslash}p{(\columnwidth - 2\tabcolsep) * \real{0.7310}}@{}}
\toprule\noalign{}
\endhead
\bottomrule\noalign{}
\endlastfoot
\begin{minipage}[t]{\linewidth}\raggedright
\begin{quote}
\textbf{Subject}
\end{quote}
\end{minipage} & Computational materials science \\
\begin{minipage}[t]{\linewidth}\raggedright
\begin{quote}
\textbf{Specific subject area}
\end{quote}
\end{minipage} & Composition-based experimental band gap prediction \\
\begin{minipage}[t]{\linewidth}\raggedright
\begin{quote}
\textbf{Type of data}
\end{quote}
\end{minipage} & Table

Figure \\
\begin{minipage}[t]{\linewidth}\raggedright
\begin{quote}
\textbf{How the data were acquired}
\end{quote}
\end{minipage} & The data was obtained by executing CrabNet v2.0.8
(available at https://github.com/sparks-baird/CrabNet) for each of the
five folds of the Matbench experimental band gap task
(https://matbench.materialsproject.org/Leaderboards\%20Per-Task/matbench\_v0.1\_matbench\_expt\_gap/).
Python code in
https://github.com/sparks-baird/matsci-opt-benchmarks/blob/7c4346624895a7826ada07ff5e44c2f49eb42b9d/scripts/crabnet\_hyperparameter/crabnet\_hyperparameter\_submitit.py
was used for orchestration, with the University of
Utah\textquotesingle s Center for High-performance Computing (CHPC)
resources. Jobs were sent to the SLURM scheduler using Submitit
(https://github.com/facebookincubator/submitit), and results were logged
in JSON format using the MongoDB Data API. The matsci-opt-benchmarks
code used for this study can be found at
https://github.com/sparks-baird/matsci-opt-benchmarks/tree/v0.2.1
(https://dx.doi.org/10.5281/zenodo.7694289). \\
\begin{minipage}[t]{\linewidth}\raggedright
\begin{quote}
\textbf{Data format}
\end{quote}
\end{minipage} & Analyzed

Filtered

Raw \\
\begin{minipage}[t]{\linewidth}\raggedright
\begin{quote}
\textbf{Description of data collection}
\end{quote}
\end{minipage} & The Ax Platform was used to perform a quasi-random
Sobol sampling of 65536 parameter combinations, varying 23
hyperparameters within a constrained search space, with 5 repeats
resulting in a total of 327680 training runs. Out of these, 173219
completed successfully, consuming 387 RTX-2080-Ti GPU days or 4614.29
CUDA core years, resulting in 41550 unique sets. To rank repeat
simulations, the "dense" method with pct=True was used in
pandas.core.groupby.GroupBy.rank. \\
\begin{minipage}[t]{\linewidth}\raggedright
\begin{quote}
\textbf{Data source location}
\end{quote}
\end{minipage} & Free-tier Shared Cluster MongoDB Atlas Database \\
\begin{minipage}[t]{\linewidth}\raggedright
\begin{quote}
\textbf{Data accessibility}
\end{quote}
\end{minipage} & \begin{minipage}[t]{\linewidth}\raggedright
\begin{quote}
Repository name: Zenodo

Data identification number: 7694268

Direct URL to data: \url{https://doi.org/10.5281/zenodo.7694268}
\end{quote}
\end{minipage} \\
\end{longtable}

\textbf{Value of the data}

\begin{itemize}
\item
  The dataset is valuable for benchmarking adaptive design methods
  applied to high-dimensional, constrained, multi-fidelity optimization
  tasks.
\item
  Practitioners of optimization in the physical sciences can leverage
  this dataset to simulate real-world materials optimization tasks, such
  as alloy discovery, and achieve improved performance.
\item
  The dataset is also useful in gaining insights into the hyperparameter
  optimization strategies used for developing compositionally restricted
  material property prediction models.
\end{itemize}

\begin{itemize}
\tightlist
\item
\item
\item
\end{itemize}

\textbf{Objective}

Industry-relevant optimization tasks in the physical sciences are often
``hierarchical, noisy, multi-fidelity\textsuperscript{1,2},
multi-objective\textsuperscript{3,4},
high-dimensional\textsuperscript{5,6}, and non-linearly correlated while
exhibiting mixed numerical and categorical variables subject to
linear\textsuperscript{7} and non-linear
constraints.''\textsuperscript{8,9} Existing benchmark
datasets\textsuperscript{10--15}, while very useful, typically are
single-objective, single-fidelity, low-dimensional, and ignore or
simplify the influence of noise. The inclusion of heteroskedastic noise
in a surrogate model enables us to establish a "Turing test" scenario
where the surrogate model is virtually identical to the actual
simulation. This approach bridges the gap between low-cost surrogate
function evaluations using benchmark datasets and costly, real-world
objective function evaluations by considering a multi-objective,
multi-fidelity, and high-dimensional task while accounting for
heteroskedastic noise.

\textbf{Data description}

The regression dataset contains hyperparameter sets (including repeats)
spanning twenty-three hyperparameter sets and their corresponding MAE,
RMSE, computational runtimes, and model size for training CrabNet.

For histogram data for the number of successful repeats see Figure 1.

For histograms of the mean absolute error, root-mean-square error,
runtime, and model size, see Figure 2, Figure 3, Figure 4, and Figure 5,
respectively.

\includegraphics[width=4.16667in,height=4.16667in]{media/image1.png}

Figure . Histogram of number of parameter groups vs. number of
successful repeats within a given group. The lowest number of repeats
for a parameter set is 1, with approximately 2.6 repeats on average.

\includegraphics[width=3.5in,height=3.5in]{media/image2.png}

Figure . Histogram of number of training runs vs. mean absolute error
using CrabNet on the Matbench experimental band gap task.

\includegraphics[width=3.5in,height=3.5in]{media/image3.png}

Figure . Histogram of number of training runs vs. root-mean-square-error
using CrabNet on the Matbench experimental band gap task.

\includegraphics[width=3.5in,height=3.5in]{media/image4.png}

Figure . Histogram of number of training runs vs. GPU runtime on an RTX
2080-Ti using CrabNet on the Matbench experimental band gap task. The
y-axis is log-scaled.

\includegraphics[width=3.5in,height=3.5in]{media/image5.png}

Figure . Histogram of number of training runs vs. model size using
CrabNet on the Matbench experimental band gap task.

\textbf{Experimental design, materials and methods}

A vast number of CrabNet models, totaling to 173219, were trained with
different hyperparameter combinations using the Ax
platform\textquotesingle s quasi-random Sobol sampling function to
generate unique parameter combinations. While there can be other uses,
this dataset is primarily intended as a multi-objective, multi-fidelity,
high-dimensional benchmark dataset for formulation-based optimization
scenarios by scaling each of the numerical parameters to the range of 0
to 1 and applying a contrived constraint that the sum of all parameters
must equal one. To realistically capture noise in the benchmark dataset,
simulations were repeated for each quasi-random parameter combination.
To improve efficiency and reduce latency, hyperparameter sets (including
repeats) were shuffled and divided into batches, then sent to a
high-performance computing environment for asynchronous evaluation.
Despite some results not being completed due to either timeout or
preemption, this trade-off was deemed reasonable for the efficiency and
completion gains.

Results were logged to a free-tier MongoDB Atlas database and then
aggregated and prepared as machine-learning-ready datasets via Python in
Jupyter notebooks. For implementation details, see
\url{https://github.com/sparks-baird/matsci-opt-benchmarks/tree/main/scripts/crabnet_hyperparameter}
and
\url{https://github.com/sparks-baird/matsci-opt-benchmarks/tree/main/notebooks/crabnet_hyperparameter}.

\textbf{Ethics statements}

There are no statements to declare.

\textbf{CRediT author statement}

\textbf{Sterling G. Baird}: Project administration, Conceptualization,
Methodology, Software, Validation, Formal analysis, Investigation, Data
Curation, Writing - Original Draft, Writing - Review \& Editing,
Visualization. \textbf{Jeet N. Parikh}: Methodology, Software, Formal
Analysis, Data Curation, Writing - Original Draft, Writing - Review \&
Editing, \textbf{Taylor D. Sparks}: Supervision, Funding acquisition

\textbf{Acknowledgments}

Funding: This work was supported by the National Science Foundation
Division of Materials Research {[}Grant number DMR-1651668{]}.

We thank Trupti Mohanty for work and discussion related to the future
use-case of this dataset as a pseudo-materials benchmark. We acknowledge
the University of Utah's Center for High Performance Computing (CHPC)
for providing computational resources. The code and manuscript
associated with a benchmark dataset for particle packing
simulations\textsuperscript{8} served as a starting point for the code
and manuscript for this work, for which there is shared language and
document structure. We acknowledge OpenAI for providing free usage of
their research tool, ChatGPT, which was used during the review and
editing process.

\textbf{Declaration of interests}

x The authors declare that they have no known competing financial
interests or personal relationships that could have appeared to
influence the work reported in this paper.

☐ The authors declare the following financial interests/personal
relationships which may be considered as potential competing interests:

\textbf{References}

(1) Ghoreishi, S. F.; Molkeri, A.; Arróyave, R.; Allaire, D.;
Srivastava, A. Efficient Use of Multiple Information Sources in Material
Design. \emph{Acta Materialia} \textbf{2019}, \emph{180}, 260--271.
https://doi.org/10.1016/j.actamat.2019.09.009.

(2) Kandasamy, K.; Vysyaraju, K. R.; Neiswanger, W.; Paria, B.; Collins,
C. R.; Schneider, J.; Poczos, B.; Xing, E. P. Tuning Hyperparameters
without Grad Students: Scalable and Robust Bayesian Optimisation with
Dragonfly. \emph{arXiv:1903.06694 {[}cs, stat{]}} \textbf{2020}.

(3) Hanaoka, K. Comparison of Conceptually Different Multi-Objective
Bayesian Optimization Methods for Material Design Problems.
\emph{Materials Today Communications} \textbf{2022}, 103440.
https://doi.org/10.1016/j.mtcomm.2022.103440.

(4) Häse, F.; Roch, L. M.; Aspuru-Guzik, A. Chimera: Enabling Hierarchy
Based Multi-Objective Optimization for Self-Driving Laboratories.
\emph{Chem. Sci.} \textbf{2018}, \emph{9} (39), 7642--7655.
https://doi.org/10.1039/C8SC02239A.

(5) Baird, S. G.; Liu, M.; Sparks, T. D. High-Dimensional Bayesian
Optimization of 23 Hyperparameters over 100 Iterations for an
Attention-Based Network to Predict Materials Property: A Case Study on
CrabNet Using Ax Platform and SAASBO. \emph{Computational Materials
Science} \textbf{2022}, \emph{211}, 111505.
https://doi.org/10.1016/j.commatsci.2022.111505.

(6) Eriksson, D.; Jankowiak, M. High-Dimensional Bayesian Optimization
with Sparse Axis-Aligned Subspaces. \emph{arXiv:2103.00349 {[}cs,
stat{]}} \textbf{2021}.

(7) Baird, S.; Hall, J. R.; Sparks, T. D. The Most Compact Search Space
Is Not Always the Most Efficient: A Case Study on Maximizing Solid
Rocket Fuel Packing Fraction via Constrained Bayesian Optimization.
ChemRxiv September 6, 2022.
https://doi.org/10.26434/chemrxiv-2022-nz2w8-v2.

(8) Baird, S. G.; Sparks, T. D. Materials Science Optimization Benchmark
Dataset for Multi-Fidelity Hard-Sphere Packing Simulations. ChemRxiv
January 9, 2023. https://doi.org/10.26434/chemrxiv-2023-fjjk7.

(9) Baird, S.; Hall, J. R.; Sparks, T. D. \emph{Compactness Matters:
Improving Bayesian Optimization Efficiency of Materials Formulations
through Invariant Search Spaces}; preprint; Chemistry, 2023.
https://doi.org/10.26434/chemrxiv-2022-nz2w8-v3.

(10) Dunn, A.; Wang, Q.; Ganose, A.; Dopp, D.; Jain, A. Benchmarking
Materials Property Prediction Methods: The Matbench Test Set and
Automatminer Reference Algorithm. \emph{npj Comput Mater} \textbf{2020},
\emph{6} (1), 138. https://doi.org/10.1038/s41524-020-00406-3.

(11) De Breuck, P.-P.; Evans, M. L.; Rignanese, G.-M. Robust Model
Benchmarking and Bias-Imbalance in Data-Driven Materials Science: A Case
Study on MODNet. \emph{J. Phys.: Condens. Matter} \textbf{2021},
\emph{33} (40), 404002. https://doi.org/10.1088/1361-648X/ac1280.

(12) Wang, A.; Liang, H.; McDannald, A.; Takeuchi, I.; Kusne, A. G.
Benchmarking Active Learning Strategies for Materials Optimization and
Discovery. arXiv April 12, 2022. http://arxiv.org/abs/2204.05838
(accessed 2022-07-04).

(13) Liang, Q.; Gongora, A. E.; Ren, Z.; Tiihonen, A.; Liu, Z.; Sun, S.;
Deneault, J. R.; Bash, D.; Mekki-Berrada, F.; Khan, S. A.;
Hippalgaonkar, K.; Maruyama, B.; Brown, K. A.; Fisher III, J.;
Buonassisi, T. Benchmarking the Performance of Bayesian Optimization
across Multiple Experimental Materials Science Domains. \emph{npj Comput
Mater} \textbf{2021}, \emph{7} (1), 188.
https://doi.org/10.1038/s41524-021-00656-9.

(14) Henderson, A. N.; Kauwe, S. K.; Sparks, T. D. Benchmark Datasets
Incorporating Diverse Tasks, Sample Sizes, Material Systems, and Data
Heterogeneity for Materials Informatics. \emph{Data in Brief}
\textbf{2021}, \emph{37}, 107262.
https://doi.org/10.1016/j.dib.2021.107262.

(15) Häse, F.; Aldeghi, M.; Hickman, R. J.; Roch, L. M.; Christensen,
M.; Liles, E.; Hein, J. E.; Aspuru-Guzik, A. Olympus: A Benchmarking
Framework for Noisy Optimization and Experiment Planning. \emph{Mach.
Learn.: Sci. Technol.} \textbf{2021}, \emph{2} (3), 035021.
https://doi.org/10.1088/2632-2153/abedc8.
